% Part: first-order-logic
% Chapter: introduction
% Section: soundness-completeness

\documentclass[../../../include/open-logic-section]{subfiles}

\begin{document}

\olfileid{fol}{int}{scp}

\olsection{నిర్దుష్టత మరియు సంపూర్ణత}

మొదటిస్థాయి విధేయ తర్కానికి \tetoken{వ్యుత్పత్తి}{!!{derivation}} వ్యవస్థలను కూడా ప్రవేశపెడతాం.
తర్కవేత్తలు అభివృద్ధి చేసిన అనేక ప్రామాణిక \tetoken{వ్యుత్పత్తి}{!!{derivation}} వ్యవస్థలు
ఉన్నాయి; కానీ అవన్నీ \tetoken{వాక్యాలు}{!!{sentence}s} మధ్య ఒకే \tetoken{వ్యుత్పాద్యత}{!!{derivability}}
సంబంధాన్ని నిర్వచిస్తాయి. ఖచ్చితంగా నిర్వచించిన ఒక నిర్దిష్ట రకపు
\tetoken{ఒక వ్యుత్పత్తి}{!!a{derivation}} ఉంటే, $\Gamma$ అనేది $!A$ను \emph{\tetoken{నిగమనం}{!!{derive}s}}
చేస్తుందని, $\Gamma \Proves !A$ అని అంటాం. \tetoken{వ్యుత్పత్తులు}{!!^{derivation}s} ఎప్పుడూ
సంకేతాల పరిమిత అమరికలే---అవి \tetoken{వాక్యాలు}{!!{sentence}s} జాబితి కావచ్చు, లేదా
మరింత సంక్లిష్టమైన నిర్మాణం కావచ్చు. ఏదైనా సమితి~$\Gamma$ నుంచి ఒక
\tetoken{వాక్యం}{!!a{sentence}} అనుగమిస్తుందా అని నిర్ణయించడానికి ఒక సాధనాన్ని
అందించడమే \tetoken{వ్యుత్పత్తి}{!!{derivation}} వ్యవస్థల ఉద్దేశం. ఆ ఉద్దేశాన్ని నెరవేర్చాలంటే,
$\Gamma \Entails !A$ కావడం, $\Gamma \Proves !A$ అయినప్పుడూ, అప్పుడు
మాత్రమే కావాలి.

$\Gamma \Proves !A$ అయ్యి, $\Gamma \Entails !A$ కాకపోతే, మన
\tetoken{వ్యుత్పత్తి}{!!{derivation}} వ్యవస్థ అతిబలమైనది; అది నిరూపించాల్సినదానికంటే ఎక్కువ
నిరూపిస్తుంది. $\Gamma \Proves !A$ అయితే $\Gamma \Entails !A$ అనే
ధర్మాన్ని \emph{నిర్దుష్టత} అంటారు; ఏ మంచి \tetoken{వ్యుత్పత్తి}{!!{derivation}} వ్యవస్థకైనా
ఇది కనీస ఆవశ్యకత. మరోవైపు, $\Gamma \Entails !A$ అయ్యి, $\Gamma
\Proves !A$ కాకపోతే, మన \tetoken{వ్యుత్పత్తి}{!!{derivation}} వ్యవస్థ చాలా బలహీనమైనది; అది
తగినంత నిరూపించదు. $\Gamma \Entails !A$ అయితే $\Gamma \Proves !A$
అనే ధర్మాన్ని \emph{సంపూర్ణత} అంటారు. నిర్దుష్టతను నిరూపించడం
సాధారణంగా సాపేక్షంగా సులభం (ఆగమనాత్మకంగా నిర్వచించిన
\tetoken{వ్యుత్పత్తులు}{!!{derivation}s} నిర్మాణంపై ఆగమనం ద్వారా). సంపూర్ణతను నిరూపించడం
మరింత కష్టం.

నిర్దుష్టతకు, సంపూర్ణతకు అనేక ముఖ్యమైన పరిణామాలు ఉన్నాయి. ఒక
\tetoken{వాక్యాలు}{!!{sentence}s} సమితి~$\Gamma$, $!A \land \lnot !A$ వంటి వైరుధ్యాన్ని
\tetoken{నిగమనం}{!!{derive}s} చేస్తే దాన్ని \emph{వైరుధ్యమైనది} అంటాం. వైరుధ్యమైన
$\Gamma$లకు నమూనాలు ఉండవు; అవి సంతృప్తిపరచలేనివి. సంపూర్ణత నుంచి
దాని విపర్యయం అనుగమిస్తుంది: వైరుధ్యం కాని---లేదా మనం అనబోయే విధంగా
\emph{అవైరుధ్యమైన}---ప్రతి $\Gamma$కు ఒక నమూనా ఉంటుంది. నిజానికి,
ఇది సంపూర్ణతకు సరిసమానం; మనం నిజంగా నిరూపించబోయే సంపూర్ణతా రూపం
ఇదే. ఇది లోతైన, బహుశా ఆశ్చర్యకరమైన ఫలితం: $\Gamma$ నుంచి
$!A \land \lnot !A$ను నిరూపించలేమనే ఒక్క వాస్తవమే, $\Gamma$
వర్ణించినట్లుగా ఉండే ఒక \tetoken{నిర్మాణం}{!!a{structure}} ఉందని హామీ ఇస్తుంది. కాబట్టి
ఏ \tetoken{వాక్యాలు}{!!{sentence}s} సమితులకు నమూనాలు ఉంటాయి అనే ప్రశ్నకు సంపూర్ణత
సమాధానం ఇస్తుంది. సమాధానం: అవైరుధ్య సమితులకు, వాటికి మాత్రమే.

నిర్దుష్టత, సంపూర్ణతా సిద్ధాంతాలకు రెండు ముఖ్యమైన పరిణామాలు ఉన్నాయి:
సంహతత్వ సిద్ధాంతం, లొవెన్‌హైమ్--స్కోలెమ్ సిద్ధాంతం. నమూనాల
సిద్ధాంతంలో ఇవి ముఖ్యమైన ఫలితాలు; అనేక ఆసక్తికరమైన ఫలితాలను
స్థాపించడానికి వీటిని ఉపయోగించవచ్చు. ఇప్పటికే రెండింటిని
ప్రస్తావించాం: ఒక \tetoken{నిర్మాణం}{!!a{structure}} యొక్క \tetoken{వ్యక్తి క్షేత్రం}{!!{domain}} పరిమితమని గానీ,
అది \tetoken{లెక్కించలేని}{!!{nonenumerable}} అని గానీ మొదటిస్థాయి విధేయ తర్కం వ్యక్తీకరించలేదు.

చారిత్రకంగా, ఈ విషయాలన్నీ---మొదటిస్థాయి విధేయ తర్కపు వాక్యనిర్మాణం,
అర్థవిచారాన్ని ఎలా నిర్వచించాలి, మంచి \tetoken{వ్యుత్పత్తి}{!!{derivation}} వ్యవస్థలను ఎలా
నిర్వచించాలి, అవి నిర్దుష్టమూ సంపూర్ణమూ అని ఎలా నిరూపించాలి,
మొదటిస్థాయి భాషల్లో దేన్ని వ్యక్తీకరించవచ్చు లేదా వ్యక్తీకరించలేమో ఎలా
స్పష్టపరచాలి---కనుగొని సరిగా రూపొందించడానికి చాలా కాలం పట్టింది.
ఇప్పుడు అది ఎలా చేయాలో మనకు తెలుసు; అయినా వివరాలన్నింటిని
అనుసరించడం గందరగోళంగా, శ్రమతో కూడినదిగా ఉండవచ్చు. కానీ అది
ముఖ్యమైనదీ కూడా; ఎందుకంటే మొదటిస్థాయి విధేయ తర్కపు ఔపచారిక భాష కోసం ఇక్కడ
అభివృద్ధి చేసిన పద్ధతులను తర్కం, కంప్యూటర్ విజ్ఞానం, భాషాశాస్త్రంలో
అన్నిచోట్లా అన్వయిస్తారు. అందువల్ల వివరాలను శ్రద్ధగా అధ్యయనం చేయడం
దీర్ఘకాలంలో ఫలిస్తుంది.

\end{document}
